% =====================================================================
%  CHAPTER 1B: 脂类 — 参考：0613笔记 (PRIMARY SOURCE)
% =====================================================================
\chapter{脂类}
\chapterintro{掌握脂类的分类和生理功能；掌握脂肪酸分类（按碳链长度/饱和程度/空间结构）、脂肪酸命名Cx:y及必需脂肪酸（亚油酸/$\alpha$-亚麻酸）；熟悉食物脂类营养价值评价（消化率/EFA含量/脂肪酸比例/脂溶性维生素含量）。}

\section{脂类的分类和生理功能}

\begin{keybox}
\textbf{脂类（Lipids）的分类：}脂类 = 脂肪（Fat）和类脂（Lipoids）。脂肪 = 甘油三酯（Triglycerides），体内重要的储能和供能物质，\imp{占体内脂类总量的95\%}。类脂：主要包括\imp{磷脂}和\imp{固醇类}，是细胞膜、组织器官，尤其是神经组织的重要组成。

\textbf{脂类的生理功能：}
\begin{enumerate}[leftmargin=*]
    \item \imp{储存和提供能量：}1g脂肪 = 9kcal能量。脂肪细胞可无限地储存脂肪。脂肪不能给脑和神经细胞以及血细胞供能。
    \item \imp{保温及润滑。}
    \item \imp{节约蛋白质作用：}保护体内蛋白质不被用作能源物质。
    \item \imp{机体构成成分：}细胞膜。
    \item \imp{脂肪组织内分泌功能：}雌激素等。
    \item \imp{食物中脂肪的作用：}增加饱腹感、改善食物感官性状、提供脂溶性维生素。
\end{enumerate}
\end{keybox}

\section{脂肪酸的分类}

\begin{keybox}
\textbf{脂肪酸命名：}表示为Cx:y，其中x为碳原子个数、y表示不饱和双键个数。

\textbf{1. 按碳链长度分类：}
\begin{itemize}[leftmargin=*]
    \item \imp{长链脂肪酸：}含14\~{}24碳
    \item \imp{中链脂肪酸（MCFA）：}含8\~{}12个碳。可直接与甘油三酯酯化；水溶性好，能直接被小肠吸收；无需形成乳糜颗粒，直接进入肝脏；细胞内快速氧化供能，但不能用于糖尿病、酮中毒及酸中毒、肝硬化者
    \item \imp{短链脂肪酸（SCFA）：}含6碳以下（醋酸、丙酸、丁酸）。来源：膳食纤维、抗性淀粉、低聚糖、糖醇等，结肠被肠道微生物发酵。功能：供能、促进细胞膜脂类物质合成、可能预防和治疗溃疡性结肠炎、可预防结肠肿瘤、对内源性胆固醇的合成有抑制作用
\end{itemize}
食物中主要以18碳脂肪酸为主，人体只能利用偶数个数碳原子的脂肪酸。

\textbf{2. 按饱和程度分类：}
\begin{itemize}[leftmargin=*]
    \item \imp{饱和脂肪酸（SFA）：}碳链中没有不饱和双键。HDL$\downarrow$，LDL$\uparrow$
    \item \imp{不饱和脂肪酸（USFA）：}碳链中含一个或多个不饱和双键
    \item \imp{单不饱和脂肪酸（MUFA）：}含一个不饱和双键，油酸。HDL$\uparrow$、LDL$\downarrow$
    \item \imp{多不饱和脂肪酸（PUFA）：}含两个及以上不饱和双键，膳食中主要为亚麻酸、$\alpha$-亚麻酸，存在于植物油中。HDL$\downarrow$、LDL$\downarrow\downarrow$
\end{itemize}

\textbf{3. 按空间结构分类：}
\begin{itemize}[leftmargin=*]
    \item \imp{顺式脂肪酸：}大多不饱和脂肪酸
    \item \imp{反式脂肪酸：}天然的不存在，主要存在于牛奶、人造奶油中。HDL$\downarrow$、LDL$\uparrow$ $\rightarrow$ 增加心血管疾病的危险。可能诱发肿瘤、2型糖尿病。人造奶油、蛋糕、饼干、油炸食品、乳酪产品、花生酱等是反式脂肪酸的主要来源。
\end{itemize}
\end{keybox}

\section{必需脂肪酸}

\begin{keybox}
\textbf{必需脂肪酸（Essential Fatty Acid, EFA）：}人体不可缺少且自身不能合成，必须通过食物供给的脂肪酸，有\imp{亚油酸}（n-6系多不饱和）和\imp{$\alpha$-亚麻酸}（n-3系多不饱和）。

\textbf{EFA的功能：}
\begin{itemize}[leftmargin=*]
    \item 构成磷脂的组成成分 $\rightarrow$ 构成细胞膜
    \item 前列腺素PG合成的前体 $\rightarrow$ 血管舒缩、神经传导等
    \item 参与胆固醇代谢
    \item 参与类十二烷酸合成（PG、TXA等），调节血压、血脂、血栓形成及免疫反应等
\end{itemize}
\end{keybox}

\section{食物脂类营养价值评价}

\begin{defbox}
\textbf{食物脂类营养价值评价：}
\begin{enumerate}[leftmargin=*]
    \item \imp{脂肪的消化率：}含不饱和脂肪酸、短链脂肪酸越多的脂肪 $\rightarrow$ 熔点越低 $\rightarrow$ 越易消化。一般植物脂肪消化率 > 动物脂肪
    \item \imp{必需脂肪酸（EFA）的含量：}一般植物油中亚油酸和$\alpha$-亚麻酸的含量 > 动物脂肪（椰子油除外）
    \item \imp{各种脂肪酸的比例：}一般推荐SFA:MUFA:PUFA = 1:1:1。EFA摄入不少于总能量的3\%，n-3系:n-6系 = 1:4\~{}6
    \item \imp{脂溶性维生素含量：}脂溶性维生素含量高的脂类 $\rightarrow$ 营养价值高。植物油（特别谷类种子的胚油）VE含量丰富。器官脂肪（肝脏脂肪）含丰富VA、VD，海产鱼类肝脏脂肪中VA、VD更高
\end{enumerate}

\textbf{参考摄入量：}成人脂肪摄入量应占总能量的\imp{20\%\~{}30\%}。EFA摄入不少于总能量的3\%。
\end{defbox}

\section{复习题}
\begin{enumerate}[leftmargin=*, itemsep=8pt]
    \item \textbf{【名词解释】}脂类（Lipids）；甘油三酯（Triglycerides）；必需脂肪酸（EFA）；脂肪酸命名Cx:y
    \item \textbf{【名词解释】}中链脂肪酸（MCFA）；短链脂肪酸（SCFA）；反式脂肪酸
    \item \textbf{【简答题】}简述脂类的分类和生理功能。
    \item \textbf{【简答题】}按碳链长度、饱和程度、空间结构分别简述脂肪酸的分类及特点。
    \item \textbf{【简答题】}简述必需脂肪酸（EFA）的定义、种类和功能。
    \item \textbf{【简答题】}简述食物脂类营养价值评价的四个方面及成人脂肪参考摄入量。
\end{enumerate}

\subsection*{参考答案要点}
\begin{enumerate}[leftmargin=*, itemsep=5pt]
    \item \textbf{脂类（Lipids）：}脂类 = 脂肪（Fat）和类脂（Lipoids）。脂肪 = 甘油三酯（Triglycerides），占体内脂类总量的95\%，是体内重要的储能和供能物质，1g脂肪 = 9kcal能量。类脂主要包括磷脂和固醇类，是细胞膜、组织器官，尤其是神经组织的重要组成。\textbf{EFA}：人体不可缺少且自身不能合成，必须通过食物供给的脂肪酸，有亚油酸（n-6系）和$\alpha$-亚麻酸（n-3系）。\textbf{Cx:y}：x为碳原子个数，y表示不饱和双键个数。
    \item \textbf{MCFA}：含8\~{}12个碳，可直接与甘油三酯酯化，水溶性好，能直接被小肠吸收，无需形成乳糜颗粒直接进入肝脏，细胞内快速氧化供能，但不能用于糖尿病、酮中毒及酸中毒、肝硬化者。\textbf{SCFA}：含6碳以下（醋酸、丙酸、丁酸），来源于膳食纤维等在结肠被肠道微生物发酵，功能包括供能、促进细胞膜脂类物质合成、可能预防和治疗溃疡性结肠炎、可预防结肠肿瘤、对内源性胆固醇的合成有抑制作用。\textbf{反式脂肪酸}：天然的不存在，主要存在于牛奶、人造奶油中，HDL$\downarrow$/LDL$\uparrow$，增加心血管疾病的危险，可能诱发肿瘤、2型糖尿病。人造奶油、蛋糕、饼干、油炸食品、乳酪产品、花生酱等是主要来源。
    \item 分类：脂类 = 脂肪（甘油三酯，占体内脂类总量的95\%） + 类脂（磷脂和固醇类，细胞膜和神经组织的重要组成）。生理功能6项：(1)储存和提供能量（1g脂肪=9kcal，脂肪细胞可无限储存脂肪，脂肪不能给脑和神经细胞以及血细胞供能）；(2)保温及润滑；(3)节约蛋白质作用；(4)机体构成成分（细胞膜）；(5)脂肪组织内分泌功能（雌激素等）；(6)食物中脂肪的作用（增加饱腹感、改善食物感官性状、提供脂溶性维生素）。
    \item \textbf{按碳链长度：}长链（14\~{}24C）、中链MCFA（8\~{}12C，水溶性好，直接吸收进入肝脏，快速氧化供能，但糖尿病/酮中毒/酸中毒/肝硬化者不能用）、短链SCFA（$\leq$6C，来自膳食纤维发酵，供能/预防结肠肿瘤/抑制内源性胆固醇合成）。\textbf{按饱和程度：}SFA（无双键，HDL$\downarrow$/LDL$\uparrow$）、MUFA（1个双键，HDL$\uparrow$/LDL$\downarrow$）、PUFA（$\geq$2个双键，HDL$\downarrow$/LDL$\downarrow\downarrow$）。\textbf{按空间结构：}顺式脂肪酸（大多不饱和脂肪酸）、反式脂肪酸（天然不存在，HDL$\downarrow$/LDL$\uparrow$，增加心血管疾病危险）。
    \item EFA定义：人体不可缺少且自身不能合成，必须通过食物供给的脂肪酸。种类：亚油酸（n-6系多不饱和）和$\alpha$-亚麻酸（n-3系多不饱和）。功能：构成磷脂的组成成分$\rightarrow$构成细胞膜；前列腺素PG合成的前体$\rightarrow$血管舒缩、神经传导等；参与胆固醇代谢；参与类十二烷酸合成（PG、TXA等），调节血压、血脂、血栓形成及免疫反应等。
    \item 四方面：(1)脂肪的消化率（含不饱和脂肪酸/短链脂肪酸越多$\rightarrow$熔点越低$\rightarrow$越易消化，植物脂肪消化率 > 动物脂肪）；(2)EFA含量（植物油 > 动物脂肪，椰子油除外）；(3)各种脂肪酸比例（SFA:MUFA:PUFA = 1:1:1，EFA不少于总能量3\%，n-3:n-6 = 1:4\~{}6）；(4)脂溶性维生素含量（植物油胚油VE丰富，器官脂肪和海产鱼类肝脏脂肪含丰富VA/VD）。成人脂肪摄入量应占总能量的20\%\~{}30\%，EFA摄入不少于总能量的3\%。
\end{enumerate}
\newpage
